\documentclass[../../../include/open-logic-section]{subfiles}

\begin{document}

\olfileid[hi]{sth}{card-arithmetic}{ch}

\olsection{सातत्य परिकल्पना}

पिछला परिणाम (सही रूप से) संकेत करता है कि यदि यह प्रत्याभूत हो कि अनंत
कार्डिनल $2$ की घातों के साथ, अर्थात्
(\olref[opps]{lem:SizePowerset2Exp} से) घात समुच्चय लेने के साथ,
``सीधे-सादे ढंग से व्यवहार'' करते हैं, तो कार्डिनल घातांक काफ़ी
\emph{सरल} होगा। किंतु हम यह नहीं मान सकते कि अनंत कार्डिनल घात समुच्चयों
के साथ वास्तव में सीधे-सादे ढंग से व्यवहार \emph{करते हैं}।

इसे खोलकर समझने के लिए हम पहले कुछ सुविधाजनक संकेत प्रस्तुत करते हैं।

\begin{defn}
जहाँ $\cardsucc{\cardfont{a}}$, $\cardfont{a}$ से पूर्णतः बड़ा लघुतम
कार्डिनल है, हम दो अनंत अनुक्रम परिभाषित करते हैं:
\begin{align*}
	\aleph_{0} &\defis \omega & 		
	\beth_{0} &\defis \omega\\
	\aleph_{\alpha \ordplus 1} &\defis \cardsucc{(\aleph_{\alpha})} &
	\beth_{\alpha+1} &\defis \cardexpo{2}{\beth_{\alpha}}\\
	\aleph_{\alpha} &\defis \bigcup_{\beta< \alpha} \aleph_{\beta} &
	\beth_{\alpha} &\defis \bigcup_{\beta < \alpha}\beth_{\beta} & \text{जब $\alpha$ सीमा क्रमसूचक हो।}
	\end{align*}
\end{defn}

$\cardsucc{\cardfont{a}}$ की परिभाषा उचित है, क्योंकि
\olref[cardinals][classing]{lem:NoLargestCardinal} हमें बताता है कि प्रत्येक
कार्डिनल $\cardfont{a}$ के लिए $\cardfont{a}$ से बड़ा कोई कार्डिनल है, और
परिमितेतर आगमन प्रत्याभूत करता है कि $\cardfont{a}$ से बड़ा कोई
\emph{लघुतम} कार्डिनल है। $\aleph$ अनुक्रम की शेष परिभाषा परिमितेतर
पुनरावर्तन से दी जाती है।

कैंटर ने यह ``$\aleph$'' संकेत प्रस्तुत किया; यह \emph{आलेफ़} है, हिब्रू
वर्णमाला का पहला अक्षर और ``अनंत'' के हिब्रू शब्द का पहला अक्षर। पियर्स ने
``$\beth$'' संकेत प्रस्तुत किया; यह \emph{बेथ} है, जो हिब्रू वर्णमाला का
दूसरा अक्षर है।\footnote{पियर्स ने दिसंबर 1900 में कैंटर को लिखे एक पत्र में
इस संकेत का उपयोग किया। दुर्भाग्य से पियर्स ने वहाँ यह त्रुटिपूर्ण तर्क भी
दिया कि $\alpha\geq\omega$ के लिए $\beth_\alpha$ का अस्तित्व नहीं है।}
अब ये संकेत हमें अनंत कार्डिनल देते हैं।

\begin{prop}
प्रत्येक क्रमसूचक $\alpha$ के लिए $\aleph_\alpha$ और $\beth_\alpha$
कार्डिनल हैं।
\end{prop}

\begin{proof}
दोनों परिणाम सरल परिमितेतर आगमन से सिद्ध होते हैं।
\olref[cardinals][classing]{omegaisacardinal} से
$\aleph_0=\beth_0=\omega$ एक कार्डिनल है। यह मानकर कि $\aleph_\alpha$ और
$\beth_\alpha$ दोनों कार्डिनल हैं, $\aleph_{\alpha+1}$ और
$\beth_{\alpha+1}$ स्पष्ट रूप से कार्डिनलों के रूप में परिभाषित हैं। और
\olref[cardinals][classing]{unioncardinalscardinal} से कार्डिनलों के किसी
समुच्चय का सम्मिलन कार्डिनल है।
\end{proof}
\noindent
इसके अतिरिक्त प्रत्येक अनंत कार्डिनल एक $\aleph$ है:

\begin{prop}
यदि $\cardfont{a}$ कोई अनंत कार्डिनल है, तो किसी अद्वितीय $\gamma$ के लिए
$\cardfont{a}=\aleph_\gamma$।
\end{prop}

\begin{proof}
कार्डिनलों पर परिमितेतर आगमन से। आगमन के लिए मान लें कि यदि
$\cardfont{b}<\cardfont{a}$, तो
$\cardfont{b}=\aleph_{\gamma_\cardfont{b}}$। यदि किसी $\cardfont{b}$ के
लिए $\cardfont{a}=\cardsucc{\cardfont{b}}$, तो
$\cardfont{a}=\cardsucc{(\aleph_{\gamma_\cardfont{b}})}=
\aleph_{\gamma_\cardfont{b}+1}$। यदि $\cardfont{a}$ किसी कार्डिनल का
उत्तरवर्ती नहीं है, तो चूँकि कार्डिनल क्रमसूचक होते हैं,
$\cardfont{a}=\bigcup_{\cardfont{b}<\cardfont{a}}\cardfont{b}=
\bigcup_{\cardfont{b}<\cardfont{a}}{\aleph_{\gamma_\cardfont{b}}}$,
अतः $\cardfont{a}=\aleph_\gamma$, जहाँ
$\gamma=\bigcup_{\cardfont{b}<\cardfont{a}}\gamma_\cardfont{b}$।
\end{proof}

चूँकि प्रत्येक अनंत कार्डिनल एक $\aleph$ है, यह हमें पूछने के लिए प्रेरित
करता है: क्या प्रत्येक अनंत कार्डिनल एक $\beth$ है? यदि यह वास्तव में
\emph{सत्य होता}, तो अनंत कार्डिनल घात समुच्चय लेने की संक्रिया के साथ
``सीधे-सादे ढंग से व्यवहार'' करते। वास्तव में तब हमारे पास निम्न होता:

\begin{defish}
\emph{सामान्यीकृत सातत्य परिकल्पना} (GCH)। सभी $\alpha$ के लिए
$\aleph_\alpha=\beth_\alpha$।
\end{defish}

इसके अतिरिक्त यदि GCH सत्य होता, तो हम कार्डिनल घातांक में कुछ पर्याप्त
सरलीकरण कर सकते थे। विशेषतः हम दिखा सकते थे कि
$\cardfont{b}<\cardfont{a}$ होने पर
$\cardexpo{\cardfont{a}}{\cardfont{b}}$ का मान
$\cardfont{a}\leq\cardexpo{\cardfont{a}}{\cardfont{b}}\leq
\cardsucc{\cardfont{a}}$ के बीच रहता है। फिर हम वे सटीक शर्तें दे सकते थे
जो निर्धारित करती हैं कि दोनों में से कौन-सी संभावना घटित होती है (अर्थात्
$\cardfont{a}=\cardexpo{\cardfont{a}}{\cardfont{b}}$ अथवा
$\cardexpo{\cardfont{a}}{\cardfont{b}}=\cardsucc{\cardfont{a}}$)।
\footnote{शर्त \emph{सह-अंतिमता} द्वारा निर्धारित होती है।}

किंतु GCH एक \emph{परिकल्पना} है, \emph{प्रमेय} नहीं। वास्तव में
\citet{Godel1938} ने सिद्ध किया कि यदि $\ZFC$ संगत है, तो
$\ZFC+\text{GCH}$ भी संगत है। किंतु बाद में पता चला कि हम उसी प्रकार
$\ZFC$ में $\lnot$GCH भी जोड़ सकते हैं। वास्तव में GCH के सबसे सरल
अतुच्छ \emph{दृष्टांत} पर विचार करें:

\begin{defish}
\emph{सातत्य परिकल्पना} (CH)। $\aleph_1=\beth_1$।
\end{defish}

\citet{Cohen1963} ने सिद्ध किया कि यदि $\ZFC$ संगत है, तो
$\ZFC+\lnot\text{CH}$ भी संगत है। अतः सातत्य परिकल्पना $\ZFC$ से स्वतंत्र
है।

इसे सातत्य परिकल्पना इसलिए कहते हैं क्योंकि ``सातत्य'' वास्तविक रेखा
$\Real$ का दूसरा नाम है। \olref[opps]{continuumis2aleph0} हमें बताता है
कि $\card{\Real}=\beth_1$। अतः सातत्य परिकल्पना कहती है कि प्राकृतिक
संख्याओं की कार्डिनलता $\aleph_0=\beth_0$ और सातत्य की कार्डिनलता
$\beth_1$ के बीच कोई कार्डिनल नहीं है।

$\ZFC$ से (G)CH की \emph{स्वतंत्रता} को देखते हुए उनकी \emph{सत्यता} के
बारे में हमें क्या कहना चाहिए? इस पर कहने के लिए \emph{बहुत कुछ} है।
वास्तव में समुच्चय सिद्धांत के हाल के बहुत-से फलदायी कार्य इन प्रश्नों की
जाँच की ओर निर्देशित रहे हैं। किंतु दो संक्षिप्त बिंदुओं पर अवश्य बल देना
चाहिए।

पहला: इन औपचारिक स्वतंत्रता परिणामों से यह \emph{तत्काल} निष्पन्न नहीं
होता कि GCH या CH में से किसी का सत्यता-मूल्य \emph{अनिर्धारित} है। संभव है
कि हमें केवल कुछ और ऐसे अभिगृहीत जोड़ने हों जो हमें स्वाभाविक लगें और जो
प्रश्न को किसी एक दिशा में तय कर दें। गेडेल ने स्वयं सुझाया कि यही उचित
प्रतिक्रिया थी।

दूसरा: $\ZFC$ से CH की स्वतंत्रता निस्संदेह \emph{चकित करने वाली} है,
किंतु वह (शाब्दिक अर्थ में) \emph{अविश्वसनीय} नहीं है। बात केवल इतनी है कि
$\ZFC$ हमें जितना बताता है, उसके अनुसार कार्डिनलों से उनके उत्तरवर्तियों
तक जाने में केवल घात समुच्चय लेने की अपेक्षा कोई अधिक सूक्ष्म साधन लग सकता
है।
 %The operation of taking powersets moves us from one stage of the
 %hierarchy to its successor stage (from $V_{\alpha}$ to
 %$V_{\alpha+1}$). 

इन दोनों प्रेक्षणों के बाद यदि आप और जानना चाहते हैं, तो अब आपको इस बहस में
हित रखने वाले विभिन्न दार्शनिकों और गणितज्ञों की ओर जाना होगा।
\footnote{यद्यपि आप \citet[\S15.6]{Potter2004} को पढ़कर आरंभ कर सकते हैं।}

\end{document}
